\documentclass[11pt]{article}
\usepackage[margin=0.9in]{geometry}
\usepackage{amsmath,amssymb}
\usepackage{xcolor}
\usepackage{textcomp}
\usepackage{fancyhdr}
\usepackage[hidelinks]{hyperref}
\definecolor{accent}{HTML}{1F6F8B}
\definecolor{soft}{HTML}{555555}
\definecolor{boxbg}{HTML}{F4F8FA}
\setlength{\parindent}{0pt}
\setlength{\parskip}{4pt}
\pagestyle{fancy}\fancyhf{}
\renewcommand{\headrulewidth}{0.4pt}
\lhead{\small\color{soft}FE Environmental \textemdash\ PEFEPrep}
\rhead{\small\color{soft}\rightmark}
\cfoot{\small\color{soft}\thepage}
\newcommand{\correct}[1]{\textbf{#1}\;{\color{accent}$\checkmark$}}
\newcommand{\optline}[2]{\par\noindent\hspace{1.2em}(#1)\hspace{0.6em}#2}

\begin{document}
\begin{center}
{\Huge\bfseries\color{accent} Groundwater, Soils \& Sediments}\\[4pt]
{\large FE Environmental \textemdash\ Day 6}\\[2pt]
{\color{soft} Study set \textbullet\ 2026-06-22 \textbullet\ 20 questions \textbullet\ every numeric answer recomputed in Python}
\end{center}
\vspace{6pt}\hrule\vspace{10pt}
\markright{Groundwater, Soils \& Sediments}
\par\noindent\colorbox{boxbg}{\parbox{\dimexpr\linewidth-2\fboxsep\relax}{%
\vspace{2pt}{\small\textbf{\color{accent}Handbook fidelity.}\ Per the NCEES FE Reference Handbook v10.6 (Civil Eng. $\rightarrow$ Groundwater, p. 297--299 and Environmental Eng. $\rightarrow$ Remediation, p. 330): Darcy's Law $Q=-KA\,dh/dx$, the specific discharge $q=-K\,dh/dx$, seepage velocity $v=q/n$, the constant-head test $K=Q/(iAt_e)$, the falling-head test $K=2.303\,(aL)/(At_e)\log_{10}(h_1/h_2)$, transmissivity $T=Kb$, Dupuit's Formula (unconfined well), and the Thiem Equation (confined well) are all printed in v10.6. The Retardation Factor $R=1+(\rho_b/n_e)K_d$, $K_d=K_{oc}f_{oc}$, and Vadose Zone Penetration $D=R_vV/A$ are in the Environmental Engineering section of v10.6. Soil index properties (void ratio, porosity, etc.) appear in the Civil Engineering --- Geotechnical section. The contaminant velocity relation $v_c=v/R$ is a direct consequence of the Retardation Factor definition and is flagged as fundamental.}\vspace{2pt}}}
\vspace{10pt}
\filbreak
\textbf{\color{accent}Q1.}\quad {\small\color{soft}[D6-Q01 \textbullet\ MCQ \textbullet\ Porosity --- definition and calculation]}\par\vspace{2pt}
A soil sample has a total volume of 100 cm\textsuperscript{3} and a void volume of 35 cm\textsuperscript{3}. The porosity n of the soil is:\par\vspace{3pt}
{\small\color{soft}Relevant:}\ $n = \dfrac{V_v}{V_T}$\par\vspace{3pt}
\optline{A}{0.25}
\optline{B}{\correct{0.35}}
\optline{C}{0.46}
\optline{D}{0.54}
\par\vspace{3pt}
\vspace{2pt}\par\noindent\colorbox{boxbg}{\parbox{\dimexpr\linewidth-2\fboxsep\relax}{%
\vspace{2pt}\textbf{Answer:} (B)\par\vspace{2pt}
{\small $n = \dfrac{V_v}{V_T} = \dfrac{35}{100} = 0.35$.}\par\vspace{2pt}
{\footnotesize\color{soft}Handbook: Geotechnical Engineering --- Soil Properties \textbullet\ Ctrl-F: Porosity}\vspace{2pt}
}}
\vspace{9pt}
\filbreak
\textbf{\color{accent}Q2.}\quad {\small\color{soft}[D6-Q02 \textbullet\ MCQ \textbullet\ Void ratio to porosity]}\par\vspace{2pt}
A sandy soil has a void ratio e = 0.65. The porosity n is most nearly:\par\vspace{3pt}
{\small\color{soft}Relevant:}\ $n = \dfrac{e}{1+e}$\par\vspace{3pt}
\optline{A}{\correct{0.39}}
\optline{B}{0.54}
\optline{C}{0.65}
\optline{D}{0.71}
\par\vspace{3pt}
\vspace{2pt}\par\noindent\colorbox{boxbg}{\parbox{\dimexpr\linewidth-2\fboxsep\relax}{%
\vspace{2pt}\textbf{Answer:} (A)\par\vspace{2pt}
{\small $n = \dfrac{e}{1+e} = \dfrac{0.65}{1.65} = 0.394 \approx 0.39$.}\par\vspace{2pt}
{\footnotesize\color{soft}Handbook: Geotechnical Engineering --- Soil Properties \textbullet\ Ctrl-F: Void Ratio; Porosity}\vspace{2pt}
}}
\vspace{9pt}
\filbreak
\textbf{\color{accent}Q3.}\quad {\small\color{soft}[D6-Q03 \textbullet\ MCQ \textbullet\ Confined vs. unconfined aquifer]}\par\vspace{2pt}
An aquifer that is bounded above by an impermeable confining layer (aquitard) so that the groundwater is under pressure greater than atmospheric is called a:\par\vspace{3pt}
\optline{A}{perched aquifer}
\optline{B}{unconfined (water-table) aquifer}
\optline{C}{\correct{confined (artesian) aquifer}}
\optline{D}{vadose zone}
\par\vspace{3pt}
\vspace{2pt}\par\noindent\colorbox{boxbg}{\parbox{\dimexpr\linewidth-2\fboxsep\relax}{%
\vspace{2pt}\textbf{Answer:} (C)\par\vspace{2pt}
{\small A confined (artesian) aquifer is overlain by an impermeable confining layer; water pressure exceeds atmospheric, so water in a tightly cased well rises above the top of the aquifer. An unconfined (water-table) aquifer has a free surface exposed to atmospheric pressure.}\par\vspace{2pt}
{\footnotesize\color{soft}Handbook: Groundwater --- Well Drawdown \textbullet\ Ctrl-F: Confined aquifer; Artesian}\vspace{2pt}
}}
\vspace{9pt}
\filbreak
\textbf{\color{accent}Q4.}\quad {\small\color{soft}[D6-Q04 \textbullet\ MCQ \textbullet\ Darcy's Law --- discharge rate (SI)]}\par\vspace{2pt}
A saturated sandy formation has a hydraulic conductivity K = 4.0 $\times$ 10⁻⁴ m/s. The hydraulic gradient across a flow cross-section of A = 2.0 m\textsuperscript{2} is i = 0.015. The Darcy discharge rate Q is most nearly:\par\vspace{3pt}
{\small\color{soft}Relevant:}\ $Q = -KA\,\dfrac{dh}{dx} = KiA$ (magnitude)\par\vspace{3pt}
\optline{A}{6.0 $\times$ 10⁻⁶ m\textsuperscript{3}/s}
\optline{B}{\correct{1.2 $\times$ 10⁻⁵ m\textsuperscript{3}/s}}
\optline{C}{2.4 $\times$ 10⁻⁵ m\textsuperscript{3}/s}
\optline{D}{3.6 $\times$ 10⁻⁵ m\textsuperscript{3}/s}
\par\vspace{3pt}
\vspace{2pt}\par\noindent\colorbox{boxbg}{\parbox{\dimexpr\linewidth-2\fboxsep\relax}{%
\vspace{2pt}\textbf{Answer:} (B)\par\vspace{2pt}
{\small $Q = KiA = (4.0\times10^{-4})(0.015)(2.0) = 1.2\times10^{-5}\ \mathrm{m^3/s}$.}\par\vspace{2pt}
{\footnotesize\color{soft}Handbook: Groundwater --- Darcy's Law \textbullet\ Ctrl-F: Darcy's Law}\vspace{2pt}
}}
\vspace{9pt}
\filbreak
\textbf{\color{accent}Q5.}\quad {\small\color{soft}[D6-Q05 \textbullet\ MCQ \textbullet\ Hydraulic gradient]}\par\vspace{2pt}
Two monitoring wells are 120 m apart along the direction of groundwater flow. Water-level elevations are 23.5 m and 21.8 m above the same datum (flow is from the high well to the low well). The hydraulic gradient i is most nearly:\par\vspace{3pt}
{\small\color{soft}Relevant:}\ $i = \dfrac{\Delta h}{\Delta x}$\par\vspace{3pt}
\optline{A}{0.0085}
\optline{B}{0.0100}
\optline{C}{\correct{0.0142}}
\optline{D}{0.0200}
\par\vspace{3pt}
\vspace{2pt}\par\noindent\colorbox{boxbg}{\parbox{\dimexpr\linewidth-2\fboxsep\relax}{%
\vspace{2pt}\textbf{Answer:} (C)\par\vspace{2pt}
{\small $i = \dfrac{\Delta h}{\Delta x} = \dfrac{23.5 - 21.8}{120} = \dfrac{1.7}{120} = 0.0142$.}\par\vspace{2pt}
{\footnotesize\color{soft}Handbook: Groundwater --- Darcy's Law \textbullet\ Ctrl-F: Hydraulic Gradient}\vspace{2pt}
}}
\vspace{9pt}
\filbreak
\textbf{\color{accent}Q6.}\quad {\small\color{soft}[D6-Q06 \textbullet\ MCQ \textbullet\ Seepage velocity vs. Darcy velocity]}\par\vspace{2pt}
An aquifer has hydraulic conductivity K = 3.0 $\times$ 10⁻⁵ m/s, hydraulic gradient i = 0.010, and effective porosity n = 0.25. The average seepage (pore-water) velocity v is most nearly:\par\vspace{3pt}
{\small\color{soft}Relevant:}\ $q = Ki$ (specific discharge / Darcy velocity) \quad $v = \dfrac{q}{n} = \dfrac{Ki}{n}$ (seepage velocity)\par\vspace{3pt}
\optline{A}{3.0 $\times$ 10⁻⁷ m/s}
\optline{B}{7.5 $\times$ 10⁻⁸ m/s}
\optline{C}{\correct{1.2 $\times$ 10⁻⁶ m/s}}
\optline{D}{1.5 $\times$ 10⁻⁵ m/s}
\par\vspace{3pt}
\vspace{2pt}\par\noindent\colorbox{boxbg}{\parbox{\dimexpr\linewidth-2\fboxsep\relax}{%
\vspace{2pt}\textbf{Answer:} (C)\par\vspace{2pt}
{\small $q = Ki = (3.0\times10^{-5})(0.010) = 3.0\times10^{-7}\ \mathrm{m/s}$.  $v = q/n = 3.0\times10^{-7}/0.25 = 1.2\times10^{-6}\ \mathrm{m/s}$. The seepage velocity is larger than the Darcy velocity because water flows only through the pore space (fraction n of the total area).}\par\vspace{2pt}
{\footnotesize\color{soft}Handbook: Groundwater --- Darcy's Law \textbullet\ Ctrl-F: Seepage Velocity; Darcy Velocity}\vspace{2pt}
}}
\vspace{9pt}
\filbreak
\textbf{\color{accent}Q7.}\quad {\small\color{soft}[D6-Q07 \textbullet\ MCQ \textbullet\ Constant-head permeability test]}\par\vspace{2pt}
A constant-head permeability test on a sand sample yields: total discharge Q = 1.5 $\times$ 10⁻⁵ m\textsuperscript{3}, elapsed time tₑ = 60 s, hydraulic gradient i = 0.80, specimen cross-sectional area A = 5 $\times$ 10⁻\textsuperscript{3} m\textsuperscript{2}. The hydraulic conductivity K is most nearly:\par\vspace{3pt}
{\small\color{soft}Relevant:}\ $K = \dfrac{Q}{i\,A\,t_e}$\par\vspace{3pt}
\optline{A}{3.1 $\times$ 10⁻⁵ m/s}
\optline{B}{\correct{6.3 $\times$ 10⁻⁵ m/s}}
\optline{C}{1.3 $\times$ 10⁻⁴ m/s}
\optline{D}{2.5 $\times$ 10⁻⁴ m/s}
\par\vspace{3pt}
\vspace{2pt}\par\noindent\colorbox{boxbg}{\parbox{\dimexpr\linewidth-2\fboxsep\relax}{%
\vspace{2pt}\textbf{Answer:} (B)\par\vspace{2pt}
{\small $K = \dfrac{Q}{iAt_e} = \dfrac{1.5\times10^{-5}}{(0.80)(5\times10^{-3})(60)} = \dfrac{1.5\times10^{-5}}{0.240} = 6.25\times10^{-5}\ \mathrm{m/s} \approx 6.3\times10^{-5}\ \mathrm{m/s}$.}\par\vspace{2pt}
{\footnotesize\color{soft}Handbook: Groundwater --- Hydraulic Conductivity/Permeability \textbullet\ Ctrl-F: Constant Head; Permeability Test}\vspace{2pt}
}}
\vspace{9pt}
\filbreak
\textbf{\color{accent}Q8.}\quad {\small\color{soft}[D6-Q08 \textbullet\ MCQ \textbullet\ Transmissivity]}\par\vspace{2pt}
A confined aquifer has a hydraulic conductivity K = 1.2 $\times$ 10⁻⁴ m/s and a saturated thickness b = 25 m. The transmissivity T is most nearly:\par\vspace{3pt}
{\small\color{soft}Relevant:}\ $T = Kb$\par\vspace{3pt}
\optline{A}{1.0 $\times$ 10⁻\textsuperscript{3} m\textsuperscript{2}/s}
\optline{B}{\correct{3.0 $\times$ 10⁻\textsuperscript{3} m\textsuperscript{2}/s}}
\optline{C}{4.8 $\times$ 10⁻\textsuperscript{3} m\textsuperscript{2}/s}
\optline{D}{1.2 $\times$ 10⁻⁴ m\textsuperscript{2}/s}
\par\vspace{3pt}
\vspace{2pt}\par\noindent\colorbox{boxbg}{\parbox{\dimexpr\linewidth-2\fboxsep\relax}{%
\vspace{2pt}\textbf{Answer:} (B)\par\vspace{2pt}
{\small $T = Kb = (1.2\times10^{-4})(25) = 3.0\times10^{-3}\ \mathrm{m^2/s}$.}\par\vspace{2pt}
{\footnotesize\color{soft}Handbook: Groundwater --- Transmissivity \textbullet\ Ctrl-F: Transmissivity}\vspace{2pt}
}}
\vspace{9pt}
\filbreak
\textbf{\color{accent}Q9.}\quad {\small\color{soft}[D6-Q09 \textbullet\ Numeric \textbullet\ Darcy discharge --- USCS]}\par\vspace{2pt}
A confined aquifer has a hydraulic conductivity K = 5 $\times$ 10⁻⁵ ft/s. Two monitoring wells spaced 600 ft apart show a head drop of 4.8 ft. The flow cross-sectional area is 1,200 ft\textsuperscript{2}. Find the Darcy discharge Q (in ft\textsuperscript{3}/s).\par\vspace{3pt}
{\small\color{soft}Relevant:}\ $i = \Delta h / \Delta x$ \quad $Q = KiA$\par\vspace{3pt}
{\small\color{soft}Numeric entry.}\par\vspace{3pt}
\vspace{2pt}\par\noindent\colorbox{boxbg}{\parbox{\dimexpr\linewidth-2\fboxsep\relax}{%
\vspace{2pt}\textbf{Answer:} $\approx$ 4.8 $\times$ 10⁻⁴ ft\textsuperscript{3}/s\par\vspace{2pt}
{\small $i = 4.8/600 = 0.0080$.  $Q = KiA = (5\times10^{-5})(0.0080)(1200) = 4.8\times10^{-4}\ \mathrm{ft^3/s}$.}\par\vspace{2pt}
{\footnotesize\color{soft}Handbook: Groundwater --- Darcy's Law \textbullet\ Ctrl-F: Darcy's Law}\vspace{2pt}
}}
\vspace{9pt}
\filbreak
\textbf{\color{accent}Q10.}\quad {\small\color{soft}[D6-Q10 \textbullet\ MCQ \textbullet\ Dupuit's Formula --- unconfined well (USCS)]}\par\vspace{2pt}
An unconfined aquifer well has reached steady-state pumping. Given K = 1.5 $\times$ 10⁻⁴ ft/s, h₁ = 22 ft (saturated height at well, r₁ = 2 ft), h₂ = 32 ft (at r₂ = 400 ft), the steady pumping rate Q is most nearly:\par\vspace{3pt}
{\small\color{soft}Relevant:}\ $Q = \dfrac{\pi K(h_2^{2}-h_1^{2})}{\ln(r_2/r_1)}$\par\vspace{3pt}
\optline{A}{0.024 cfs}
\optline{B}{0.035 cfs}
\optline{C}{\correct{0.048 cfs}}
\optline{D}{0.072 cfs}
\par\vspace{3pt}
\vspace{2pt}\par\noindent\colorbox{boxbg}{\parbox{\dimexpr\linewidth-2\fboxsep\relax}{%
\vspace{2pt}\textbf{Answer:} (C)\par\vspace{2pt}
{\small $h_2^2 - h_1^2 = 32^2 - 22^2 = 1024 - 484 = 540\ \mathrm{ft^2}$.  $\ln(400/2) = \ln(200) = 5.298$.  $Q = \dfrac{\pi(1.5\times10^{-4})(540)}{5.298} = \dfrac{0.2545}{5.298} = 0.0480\ \mathrm{cfs}$.}\par\vspace{2pt}
{\footnotesize\color{soft}Handbook: Groundwater --- Dupuit's Formula \textbullet\ Ctrl-F: Dupuit; Unconfined Well}\vspace{2pt}
}}
\vspace{9pt}
\filbreak
\textbf{\color{accent}Q11.}\quad {\small\color{soft}[D6-Q11 \textbullet\ MCQ \textbullet\ Thiem Equation --- confined well (SI)]}\par\vspace{2pt}
A confined aquifer (b = 30 m, K = 8 $\times$ 10⁻⁵ m/s) is pumped to steady state. Piezometric heads: h₁ = 18 m at r₁ = 0.2 m and h₂ = 25 m at r₂ = 200 m. The pumping rate Q is most nearly:\par\vspace{3pt}
{\small\color{soft}Relevant:}\ $T = Kb$ \quad $Q = \dfrac{2\pi T(h_2-h_1)}{\ln(r_2/r_1)}$\par\vspace{3pt}
\optline{A}{7.5 $\times$ 10⁻\textsuperscript{3} m\textsuperscript{3}/s}
\optline{B}{\correct{1.5 $\times$ 10⁻\textsuperscript{2} m\textsuperscript{3}/s}}
\optline{C}{2.3 $\times$ 10⁻\textsuperscript{2} m\textsuperscript{3}/s}
\optline{D}{3.1 $\times$ 10⁻\textsuperscript{2} m\textsuperscript{3}/s}
\par\vspace{3pt}
\vspace{2pt}\par\noindent\colorbox{boxbg}{\parbox{\dimexpr\linewidth-2\fboxsep\relax}{%
\vspace{2pt}\textbf{Answer:} (B)\par\vspace{2pt}
{\small $T = Kb = (8\times10^{-5})(30) = 2.4\times10^{-3}\ \mathrm{m^2/s}$.  $\ln(200/0.2) = \ln(1000) = 6.908$.  $Q = \dfrac{2\pi(2.4\times10^{-3})(25-18)}{6.908} = \dfrac{2\pi(0.0168)}{6.908} = \dfrac{0.1056}{6.908} = 0.01528\ \mathrm{m^3/s} \approx 1.5\times10^{-2}\ \mathrm{m^3/s}$.}\par\vspace{2pt}
{\footnotesize\color{soft}Handbook: Groundwater --- Thiem Equation \textbullet\ Ctrl-F: Thiem; Confined Well}\vspace{2pt}
}}
\vspace{9pt}
\filbreak
\textbf{\color{accent}Q12.}\quad {\small\color{soft}[D6-Q12 \textbullet\ MCQ \textbullet\ Specific capacity of a well]}\par\vspace{2pt}
A municipal well is pumped at a rate Q = 0.050 cfs and the resulting drawdown at the well is Dw = 8.0 ft. The specific capacity of the well is:\par\vspace{3pt}
{\small\color{soft}Relevant:}\ $SC = \dfrac{Q}{D_w}$\par\vspace{3pt}
\optline{A}{0.0031 cfs/ft}
\optline{B}{\correct{0.0063 cfs/ft}}
\optline{C}{0.0125 cfs/ft}
\optline{D}{0.40 cfs/ft}
\par\vspace{3pt}
\vspace{2pt}\par\noindent\colorbox{boxbg}{\parbox{\dimexpr\linewidth-2\fboxsep\relax}{%
\vspace{2pt}\textbf{Answer:} (B)\par\vspace{2pt}
{\small $SC = Q/D_w = 0.050/8.0 = 0.00625\ \mathrm{cfs/ft} \approx 0.0063\ \mathrm{cfs/ft}$.}\par\vspace{2pt}
{\footnotesize\color{soft}Handbook: Groundwater --- Dupuit's Formula (specific capacity) \textbullet\ Ctrl-F: Specific Capacity; Drawdown}\vspace{2pt}
}}
\vspace{9pt}
\filbreak
\textbf{\color{accent}Q13.}\quad {\small\color{soft}[D6-Q13 \textbullet\ MCQ \textbullet\ Falling-head permeability test]}\par\vspace{2pt}
A falling-head permeability test on a clay sample gives: standpipe area a = 1 cm\textsuperscript{2}, specimen cross-sectional area A = 50 cm\textsuperscript{2}, specimen length L = 15 cm, elapsed time tₑ = 480 s, initial head h₁ = 80 cm, final head h₂ = 40 cm. The hydraulic conductivity K is most nearly:\par\vspace{3pt}
{\small\color{soft}Relevant:}\ $K = 2.303\,\dfrac{aL}{A\,t_e}\,\log_{10}\!\left(\dfrac{h_1}{h_2}\right)$\par\vspace{3pt}
\optline{A}{2.2 $\times$ 10⁻⁶ m/s}
\optline{B}{\correct{4.3 $\times$ 10⁻⁶ m/s}}
\optline{C}{8.7 $\times$ 10⁻⁶ m/s}
\optline{D}{1.4 $\times$ 10⁻⁵ m/s}
\par\vspace{3pt}
\vspace{2pt}\par\noindent\colorbox{boxbg}{\parbox{\dimexpr\linewidth-2\fboxsep\relax}{%
\vspace{2pt}\textbf{Answer:} (B)\par\vspace{2pt}
{\small $K = 2.303\,\dfrac{(1)(15)}{(50)(480)}\,\log_{10}(80/40) = 2.303\times6.25\times10^{-4}\times0.3010 = 4.33\times10^{-4}\ \mathrm{cm/s} = 4.3\times10^{-6}\ \mathrm{m/s}$.}\par\vspace{2pt}
{\footnotesize\color{soft}Handbook: Groundwater --- Hydraulic Conductivity/Permeability \textbullet\ Ctrl-F: Falling Head; Permeability Test}\vspace{2pt}
}}
\vspace{9pt}
\filbreak
\textbf{\color{accent}Q14.}\quad {\small\color{soft}[D6-Q14 \textbullet\ Numeric \textbullet\ Dupuit's Formula --- pumping rate in gpm (USCS)]}\par\vspace{2pt}
An unconfined aquifer well reaches steady state with the following data: K = 2.0 $\times$ 10⁻⁴ ft/s, well radius r₁ = 1.5 ft with saturated height h₁ = 30 ft, observation well at r₂ = 600 ft with h₂ = 45 ft. Find the steady-state pumping rate Q in gallons per minute (1 ft\textsuperscript{3}/s = 448.8 gpm).\par\vspace{3pt}
{\small\color{soft}Relevant:}\ $Q = \dfrac{\pi K(h_2^{2}-h_1^{2})}{\ln(r_2/r_1)}$\par\vspace{3pt}
{\small\color{soft}Numeric entry.}\par\vspace{3pt}
\vspace{2pt}\par\noindent\colorbox{boxbg}{\parbox{\dimexpr\linewidth-2\fboxsep\relax}{%
\vspace{2pt}\textbf{Answer:} $\approx$ 53 gpm\par\vspace{2pt}
{\small $h_2^2 - h_1^2 = 45^2 - 30^2 = 2025 - 900 = 1125\ \mathrm{ft^2}$.  $\ln(600/1.5) = \ln(400) = 5.991$.  $Q = \dfrac{\pi(2.0\times10^{-4})(1125)}{5.991} = \dfrac{0.7069}{5.991} = 0.1180\ \mathrm{ft^3/s}$.  $0.1180\times448.8 = 52.9 \approx 53\ \mathrm{gpm}$.}\par\vspace{2pt}
{\footnotesize\color{soft}Handbook: Groundwater --- Dupuit's Formula \textbullet\ Ctrl-F: Dupuit; Unconfined Well}\vspace{2pt}
}}
\vspace{9pt}
\filbreak
\textbf{\color{accent}Q15.}\quad {\small\color{soft}[D6-Q15 \textbullet\ Numeric \textbullet\ Thiem Equation --- pumping rate (SI)]}\par\vspace{2pt}
A confined aquifer has transmissivity T = 3.5 $\times$ 10⁻\textsuperscript{3} m\textsuperscript{2}/s. At steady state, piezometric heads are h₁ = 12 m at r₁ = 0.25 m (well face) and h₂ = 18 m at r₂ = 250 m. Determine the pumping rate Q (in m\textsuperscript{3}/s).\par\vspace{3pt}
{\small\color{soft}Relevant:}\ $Q = \dfrac{2\pi T(h_2-h_1)}{\ln(r_2/r_1)}$\par\vspace{3pt}
{\small\color{soft}Numeric entry.}\par\vspace{3pt}
\vspace{2pt}\par\noindent\colorbox{boxbg}{\parbox{\dimexpr\linewidth-2\fboxsep\relax}{%
\vspace{2pt}\textbf{Answer:} $\approx$ 0.0191 m\textsuperscript{3}/s\par\vspace{2pt}
{\small $\ln(250/0.25) = \ln(1000) = 6.908$.  $Q = \dfrac{2\pi(3.5\times10^{-3})(18-12)}{6.908} = \dfrac{2\pi(0.021)}{6.908} = \dfrac{0.13195}{6.908} = 0.01910\ \mathrm{m^3/s}$.}\par\vspace{2pt}
{\footnotesize\color{soft}Handbook: Groundwater --- Thiem Equation \textbullet\ Ctrl-F: Thiem; Confined Well}\vspace{2pt}
}}
\vspace{9pt}
\filbreak
\textbf{\color{accent}Q16.}\quad {\small\color{soft}[D6-Q16 \textbullet\ MCQ \textbullet\ Retardation factor --- contaminant velocity]}\par\vspace{2pt}
A contaminant in an aquifer has a retardation factor R = 4.0. If the average seepage velocity of the groundwater is 0.50 m/day, the velocity of the contaminant plume is most nearly:\par\vspace{3pt}
{\small\color{soft}Relevant:}\ $v_c = \dfrac{v}{R}$  (follows directly from $R = 1 + (\rho_b/n_e)K_d$)\par\vspace{3pt}
\optline{A}{\correct{0.125 m/day}}
\optline{B}{0.250 m/day}
\optline{C}{0.500 m/day}
\optline{D}{2.00 m/day}
\par\vspace{3pt}
\vspace{2pt}\par\noindent\colorbox{boxbg}{\parbox{\dimexpr\linewidth-2\fboxsep\relax}{%
\vspace{2pt}\textbf{Answer:} (A)\par\vspace{2pt}
{\small $v_c = v/R = 0.50/4.0 = 0.125\ \mathrm{m/day}$. The retardation factor R expresses how much slower a sorbing contaminant travels relative to the average pore-water velocity. R = 1 means no retardation (conservative tracer).}\par\vspace{2pt}
{\footnotesize\color{soft}Handbook: Remediation --- Retardation Factor \textbullet\ Ctrl-F: Retardation Factor}\vspace{2pt}
}}
\vspace{9pt}
\filbreak
\textbf{\color{accent}Q17.}\quad {\small\color{soft}[D6-Q17 \textbullet\ MCQ \textbullet\ Distribution coefficient Kd from Koc and foc]}\par\vspace{2pt}
A soil has a fraction of organic carbon foc = 0.02 and an organic carbon partition coefficient Koc = 150 mL/g. The distribution coefficient Kd is most nearly:\par\vspace{3pt}
{\small\color{soft}Relevant:}\ $K_d = K_{oc}\,f_{oc}$\par\vspace{3pt}
\optline{A}{0.013 mL/g}
\optline{B}{0.30 mL/g}
\optline{C}{\correct{3.0 mL/g}}
\optline{D}{30 mL/g}
\par\vspace{3pt}
\vspace{2pt}\par\noindent\colorbox{boxbg}{\parbox{\dimexpr\linewidth-2\fboxsep\relax}{%
\vspace{2pt}\textbf{Answer:} (C)\par\vspace{2pt}
{\small $K_d = K_{oc}\,f_{oc} = 150\times0.02 = 3.0\ \mathrm{mL/g}$.}\par\vspace{2pt}
{\footnotesize\color{soft}Handbook: Remediation --- Soil-Water Partition Coefficient / Retardation Factor \textbullet\ Ctrl-F: Koc; Partition Coefficient}\vspace{2pt}
}}
\vspace{9pt}
\filbreak
\textbf{\color{accent}Q18.}\quad {\small\color{soft}[D6-Q18 \textbullet\ MCQ \textbullet\ Vadose zone penetration depth]}\par\vspace{2pt}
A 0.5 m\textsuperscript{3} gasoline spill occurs over a surface area of 4 m\textsuperscript{2} at a site with medium to fine sand (Rv = 80). Using the Handbook vadose zone penetration formula, the maximum depth of penetration D is most nearly:\par\vspace{3pt}
{\small\color{soft}Relevant:}\ $D = \dfrac{R_v\,V}{A}$\par\vspace{3pt}
\optline{A}{5 m}
\optline{B}{\correct{10 m}}
\optline{C}{20 m}
\optline{D}{40 m}
\par\vspace{3pt}
\vspace{2pt}\par\noindent\colorbox{boxbg}{\parbox{\dimexpr\linewidth-2\fboxsep\relax}{%
\vspace{2pt}\textbf{Answer:} (B)\par\vspace{2pt}
{\small $D = \dfrac{R_v V}{A} = \dfrac{80\times0.5}{4} = \dfrac{40}{4} = 10\ \mathrm{m}$.}\par\vspace{2pt}
{\footnotesize\color{soft}Handbook: Remediation --- Vadose Zone Penetration \textbullet\ Ctrl-F: Vadose Zone; Penetration}\vspace{2pt}
}}
\vspace{9pt}
\filbreak
\textbf{\color{accent}Q19.}\quad {\small\color{soft}[D6-Q19 \textbullet\ Numeric \textbullet\ Retardation factor --- full calculation]}\par\vspace{2pt}
An aquifer has a bulk density $\rho$b = 1.8 g/cm\textsuperscript{3}, effective porosity ne = 0.30, organic carbon partition coefficient Koc = 200 cm\textsuperscript{3}/g, and soil organic carbon fraction foc = 0.01. Compute the retardation factor R.\par\vspace{3pt}
{\small\color{soft}Relevant:}\ $K_d = K_{oc}\,f_{oc}$ \quad $R = 1 + \dfrac{\rho_b}{n_e}\,K_d$\par\vspace{3pt}
{\small\color{soft}Numeric entry.}\par\vspace{3pt}
\vspace{2pt}\par\noindent\colorbox{boxbg}{\parbox{\dimexpr\linewidth-2\fboxsep\relax}{%
\vspace{2pt}\textbf{Answer:} R = 13\par\vspace{2pt}
{\small $K_d = K_{oc}\,f_{oc} = 200\times0.01 = 2.0\ \mathrm{cm^3/g}$.  $R = 1 + \dfrac{\rho_b}{n_e}K_d = 1 + \dfrac{1.8}{0.30}\times2.0 = 1 + 6.0\times2.0 = 1 + 12 = 13$. This means the contaminant travels at 1/13 the groundwater seepage velocity.}\par\vspace{2pt}
{\footnotesize\color{soft}Handbook: Remediation --- Retardation Factor \textbullet\ Ctrl-F: Retardation Factor; Koc}\vspace{2pt}
}}
\vspace{9pt}
\filbreak
\textbf{\color{accent}Q20.}\quad {\small\color{soft}[D6-Q20 \textbullet\ MCQ \textbullet\ Effective porosity --- physical meaning]}\par\vspace{2pt}
In the groundwater seepage velocity equation v = q/n, the value n is the *effective* porosity rather than total porosity because it:\par\vspace{3pt}
{\small\color{soft}Relevant:}\ $v = \dfrac{q}{n} = \dfrac{Ki}{n}$\par\vspace{3pt}
\optline{A}{includes all pore space, even that occupied by bound or trapped water}
\optline{B}{\correct{excludes pore space occupied by immobile (bound/trapped) water that does not transmit flow}}
\optline{C}{is always larger than total porosity in fine-grained soils}
\optline{D}{equals the void ratio e for all soil types}
\par\vspace{3pt}
\vspace{2pt}\par\noindent\colorbox{boxbg}{\parbox{\dimexpr\linewidth-2\fboxsep\relax}{%
\vspace{2pt}\textbf{Answer:} (B)\par\vspace{2pt}
{\small Effective porosity is the fraction of pore space through which water actually flows --- it excludes isolated pores, dead-end pores, and water held by surface tension or adsorption. Because effective porosity < total porosity, the seepage velocity v = q/n is *higher* than the Darcy velocity q; it better represents the true speed at which dissolved contaminants (non-retarded) travel.}\par\vspace{2pt}
{\footnotesize\color{soft}Handbook: Groundwater --- Darcy's Law \textbullet\ Ctrl-F: Effective Porosity; Seepage Velocity}\vspace{2pt}
}}
\vspace{9pt}
\end{document}